\chapter{Numbers up to 10\,000}\label{ch:g3:numbers}

One thousand, two thousand, three thousand\,\dots\ big numbers are
built from small ones, using a wonderful idea: the \emph{place} of a
digit tells its value. This chapter teaches how to read, write,
compare and order numbers up to $10\,000$.

\section{Reading and writing numbers}

\begin{definition}[Place value]\label{def:g3:numbers:place}
A number is written with the digits $0, 1, 2, \dots, 9$. Reading from
the right, the places are: units, tens, hundreds, thousands. In
$3\,254$:
\[
3\,254 = 3 \text{ thousands} + 2 \text{ hundreds} + 5 \text{ tens}
+ 4 \text{ units}.
\]
\end{definition}

\begin{omfigure}
\begin{tikzpicture}[scale=0.9]
  \draw (0,0) grid[xstep=2.3, ystep=0.85] (9.2,1.7);
  \node[font=\small] at (1.15,1.275) {thousands};
  \node[font=\small] at (3.45,1.275) {hundreds};
  \node[font=\small] at (5.75,1.275) {tens};
  \node[font=\small] at (8.05,1.275) {units};
  \node[font=\large, omDef] at (1.15,0.425) {$3$};
  \node[font=\large, omDef] at (3.45,0.425) {$2$};
  \node[font=\large, omDef] at (5.75,0.425) {$5$};
  \node[font=\large, omDef] at (8.05,0.425) {$4$};
\end{tikzpicture}

{\small The place-value table of $3\,254$, read ``three thousand two
hundred fifty-four''.}
\end{omfigure}

\begin{example}[The zero that holds a place]\label{ex:g3:numbers:zero}
``Four thousand and seven'' is written $4\,007$: a $4$ in the
thousands, a $7$ in the units, and \emph{zeros} in between to keep
each digit in its right place. Without them, $47$ would be a very
different number!
\end{example}

\begin{example}[Exchanging]\label{ex:g3:numbers:exchange}
Ten units make one ten; ten tens make one hundred; ten hundreds make
one thousand. So $10$ hundreds $= 1\,000$, and $32$ tens
$= 320$: the number $320$ contains $32$ tens (not just the digit
$2$!).
\end{example}

\section{Comparing and ordering}

\begin{method}[Comparing two numbers]\label{met:g3:numbers:compare}
\begin{enumerate}
  \item Count the digits: more digits means a bigger number
        ($1\,002 > 998$).
  \item Same number of digits: compare digit by digit starting from
        the \emph{left}. The first difference decides:
        $5\,614 > 5\,589$ because $6 > 5$ at the hundreds.
  \item Write the answer with the signs $<$ (smaller than) or $>$
        (bigger than): the open side always faces the bigger number.
\end{enumerate}
\end{method}

\begin{omfigure}
\begin{tikzpicture}[scale=1.1]
  \draw[-Stealth] (-0.3,0) -- (10.5,0);
  \foreach \i in {0,1,...,10} \draw (\i,-0.08) -- (\i,0.08);
  \node[below, font=\small] at (0,-0.1) {$0$};
  \node[below, font=\small] at (5,-0.1) {$500$};
  \node[below, font=\small] at (10,-0.1) {$1000$};
  \fill[omThm] (2.4,0) circle (2pt) node[above, font=\small] {$240$};
  \fill[omDef] (6.9,0) circle (2pt) node[above, font=\small] {$690$};
  \fill[omMeth] (9.05,0) circle (2pt) node[above, font=\small] {$905$};
\end{tikzpicture}

{\small Numbers on a line graduated in hundreds: further right means
bigger. $240 < 690 < 905$.}
\end{omfigure}

\begin{example}\label{ex:g3:numbers:order}
Order from smallest to biggest: $780$, $8\,700$, $807$, $87$.
First by the number of digits: $87$ (two), then $780$ and $807$
(three), then $8\,700$ (four). Between $780$ and $807$: compare the
hundreds digits, $7 < 8$, so $780 < 807$. Final order:
\[
87 < 780 < 807 < 8\,700 .
\]
\end{example}

\section{Even and odd}

\begin{definition}[Even, odd]\label{def:g3:numbers:evenodd}
A number is \emph{even}\index{even number} when it can be split into
two equal whole parts --- its units digit is $0$, $2$, $4$, $6$ or
$8$. Otherwise it is \emph{odd}\index{odd number} (units digit $1$,
$3$, $5$, $7$ or $9$).
\end{definition}

\begin{example}
$46$ is even ($46 = 23 + 23$); $47$ is odd (two equal parts would
need half an object). Only the \emph{last} digit matters: $3\,578$ is
even, $2\,341$ is odd.
\end{example}

\begin{omfigure}
\begin{tikzpicture}[scale=0.45]
  \foreach \i in {0,...,5} \fill[omDef] ({\i},0) circle (0.28);
  \foreach \i in {0,...,5} \fill[omDef] ({\i},1) circle (0.28);
  \node[below, font=\small] at (2.5,-0.7) {$12$: even, two equal rows};
  \begin{scope}[xshift=10.5cm]
  \foreach \i in {0,...,4} \fill[omThm] ({\i},0) circle (0.28);
  \foreach \i in {0,...,4} \fill[omThm] ({\i},1) circle (0.28);
  \fill[omThm] (5,0) circle (0.28);
  \draw[omExo, thick] (5,0) circle (0.5);
  \node[below, font=\small] at (2.5,-0.7) {$11$: odd, one dot alone};
  \end{scope}
\end{tikzpicture}

{\small Even numbers split into two equal rows; odd numbers always
leave one alone.}
\end{omfigure}

\section{Exercises}

\begin{exercise}[$\star$]\label{exo:g3:numbers:1}
Write in digits: ``two thousand three hundred forty-five''; ``six
thousand and fifty''; ``nine thousand and one''.
\end{exercise}

\begin{exercise}[$\star$]\label{exo:g3:numbers:2}
Write in words: $1\,208$; \quad $4\,070$; \quad $9\,999$.
\end{exercise}

\begin{exercise}[$\star$]\label{exo:g3:numbers:3}
In $7\,382$: which digit is in the tens place? In the thousands
place? How many \emph{hundreds} does the number contain in total
(careful: not just the hundreds digit --- think of
\cref{ex:g3:numbers:exchange})?
\end{exercise}

\begin{exercise}[$\star$]\label{exo:g3:numbers:4}
Decompose as in \cref{def:g3:numbers:place}:
$5\,631$; \quad $2\,046$; \quad $8\,500$.
\end{exercise}

\begin{exercise}[$\star$]\label{exo:g3:numbers:5}
Copy and complete with $<$ or $>$:
\[
456 \;?\; 465, \qquad
1\,203 \;?\; 989, \qquad
6\,090 \;?\; 6\,009, \qquad
9\,999 \;?\; 10\,000 .
\]
\end{exercise}

\begin{exercise}[$\star$]\label{exo:g3:numbers:6}
Order from smallest to biggest: $530$; \ $3\,500$; \ $353$; \ $5\,033$;
\ $503$.
\end{exercise}

\begin{exercise}[$\star$]\label{exo:g3:numbers:7}
On a number line graduated in hundreds from $0$ to $1\,000$, place
(approximately) the numbers $150$, $480$, $820$ and $990$.
\end{exercise}

\begin{exercise}[$\star$]\label{exo:g3:numbers:8}
Even or odd? \ $34$; \ $57$; \ $70$; \ $2\,481$; \ $6\,548$.
\end{exercise}

\begin{exercise}[$\star$]\label{exo:g3:numbers:9}
Count by steps: from $2\,970$, count five steps of $10$; from $860$,
count four steps of $100$; from $9\,996$, count four steps of $1$.
What do you notice at the end of the last one?
\end{exercise}

\begin{exercise}[$\star\star$]\label{exo:g3:numbers:10}
Using each of the digits $3$, $7$, $0$, $5$ exactly once, write the
biggest possible four-digit number, then the smallest one (a number
cannot start with $0$).
\end{exercise}

\begin{exercise}[$\star\star$]\label{exo:g3:numbers:11}
I am a three-digit number. My digits are $2$, $5$ and $8$, I am even,
and I am bigger than $800$. Who am I?
\end{exercise}
